Title
Probabilistic Electricity Demand Forecasting During Extreme Winter Events: A Quantile Regression Case Study of the 2014 Polar Vortex in PJM

Abstract
Accurate short-term load forecasting during extreme cold-weather events is critical for grid reliability, yet standard probabilistic evaluation frameworks aggregate performance across all operating conditions and may conceal systematic undercoverage precisely when grid stress is highest. This study evaluates gradient boosted quantile regression (QR-GBT) on the January 6–8, 2014 PJM Polar Vortex — a near-annual-peak cold-weather stress event in which RTO load reached 140,510.2 MW — using PJM hourly load data from 2010 through 2014 and same-hour ERA5 retrospective reanalysis weather inputs, which provide an idealized retrospective weather-information setting rather than an operational forecast-weather setting. On point accuracy, a gradient boosted model achieves a full-year mean absolute error of 721 MW against a PJM Day-Ahead benchmark of 1,937 MW; during the vortex window the same model records 1,705 MW MAE against 3,148 MW for the Day-Ahead benchmark, with both comparisons reflecting the ERA5 retrospective setting rather than operational superiority. Despite 86.8% empirical coverage at the nominal 90% prediction interval level across full-year 2014 evaluation, QR-GBT coverage collapses to 66.7% during the 72-hour vortex window, 98% interval coverage falls to 84.7%, and the observed load exceeds the post-rearrangement q99 estimate of 139,585 MW, leaving the event peak outside the nominal 98% prediction interval. These results demonstrate that full-year calibration metrics can mask severe tail undercoverage under distribution-shifting cold-weather stress, and that disaggregated evaluation on historical extreme windows is a necessary complement to annual probabilistic benchmarking.
Word count: 218

1. Introduction
1.1 Load Forecasting and Grid Reliability
Reliable short-term load forecasting is foundational to the secure and economical operation of bulk electric power systems. System operators depend on accurate hourly demand forecasts to schedule generating units, procure operating reserves, clear day-ahead energy markets, and manage transmission constraints [CITATION: power systems operations reference]. When forecasts underestimate demand — particularly at system peaks — the consequences are immediate: reserve margins erode, emergency procedures may be triggered, and real-time prices can spike sharply as operators scramble to secure additional capacity [CITATION: NERC reliability standards reference]. Conversely, systematic overestimation wastes reserve capacity and increases costs for consumers. The stakes are not symmetric: in a high-load emergency, a forecasting error of several gigawatts can be the difference between controlled load management and an uncontrolled reliability event. This operational asymmetry motivates particular attention to the accuracy of forecasts during extreme demand periods.
1.2 Probabilistic Forecasting and Prediction Intervals
Point forecasts — single-valued predictions of expected demand — have long been the operational standard, but they provide no direct information about uncertainty. A point forecast of 130,000 MW carries the same face value whether the forecaster is highly confident or deeply uncertain. Probabilistic load forecasting addresses this by estimating the full conditional distribution of future demand, typically expressed as quantile forecasts or prediction intervals [CITATION: probabilistic load forecasting review]. Prediction intervals support a range of reliability-critical decisions that point forecasts cannot: procuring reserves calibrated to the tail of the demand distribution, setting offer caps in capacity markets, and communicating forecast uncertainty to transmission planners [CITATION: probabilistic forecasting in power systems]. Despite these advantages, probabilistic forecasting methods are less uniformly adopted in operational practice than their deterministic counterparts, and the evaluation frameworks used to benchmark them — typically aggregate annual metrics such as the Continuous Ranked Probability Score (CRPS) or empirical prediction interval coverage — may not capture the behavior that matters most for grid reliability: performance during the small fraction of hours when load is at or near its annual extreme.
1.3 Extreme Cold as a Tail-Risk Forecasting Problem
Extreme cold-weather events present a qualitatively distinct forecasting challenge. In large interconnections with significant heating load, cold air outbreaks drive demand responses that differ materially from the moderate-winter conditions that dominate any multi-year training dataset. The load-temperature relationship is nonlinear: each additional degree of cold below the heating base temperature drives progressively larger demand increases as less-efficient heating systems engage and as behavioral responses compound meteorological forcing [CITATION: load-temperature nonlinearity reference]. A data-driven forecasting model trained primarily on moderate winters will have seen few or no examples of the load levels that arise during a severe cold air outbreak, placing the event in or near the extrapolation regime of the model's learned input-output mapping. Under these conditions, a model may produce point forecasts with acceptable average error while simultaneously constructing prediction intervals that are far too narrow — assigning insufficient probability mass to the load levels that actually occur at the peak hour. This distinction between average-case and tail-case probabilistic performance is not captured by full-year aggregate evaluation metrics, and the gap between them has direct implications for how probabilistic forecasters should be validated before being relied upon for grid reliability decisions [CITATION: distributional robustness / forecasting under distribution shift].
1.4 The January 2014 PJM Polar Vortex as a Case Study
The January 2014 polar vortex provides a well-documented and historically significant test of load forecasting under extreme cold stress. The January 6–8, 2014 Polar Vortex produced a near-annual-peak RTO load of 140,510.2 MW at Jan 7 18:00 EPT, equal to 99.18% of the 2014 annual peak of 141,677.9 MW recorded on Jun 17 17:00 EPT. This event was the subject of detailed post-event reviews by both NERC and PJM, which documented widespread generation unit failures, emergency energy transactions, and significant forecast errors by multiple market participants [CITATION: NERC Polar Vortex Review 2014; CITATION: PJM Polar Vortex Review 2014]. Several factors make this event particularly suited to a rigorous retrospective forecasting study. First, the event load is near the 2014 annual system peak, making it a genuine extreme in the demand distribution rather than a moderate stress event. Second, complete, publicly accessible hourly load data from PJM Data Miner 2 and ERA5 retrospective reanalysis weather data from the Copernicus Climate Data Store are available for exact replication of the analysis. Third, the event occurred only one year into the evaluation period, meaning a model trained through 2013 had no exposure to a comparable cold air outbreak during training — creating a clean test of out-of-distribution generalization. Fourth, the meteorological character of the event is well-characterized in the atmospheric science literature [CITATION: meteorological analysis of 2014 polar vortex], providing context for the thermodynamic forcing that drove the load response.
1.5 This Study: What Is Evaluated and What Is Found
This study evaluates a gradient boosted quantile regression model (QR-GBT) trained on 2010–2013 PJM hourly load data with ERA5 retrospective reanalysis weather features, assessed against two evaluation windows: full-year 2014 and the 72-hour January 6–8 polar vortex window. All weather-informed model results use same-hour ERA5 retrospective reanalysis weather input — that is, the actual meteorological conditions as reconstructed by the ERA5 reanalysis system rather than a real-time operational weather forecast. This design tests whether a well-specified probabilistic model can adequately characterize demand uncertainty even under an idealized weather-knowledge setting; it does not simulate operational deployment. On point accuracy, the gradient boosted model substantially outperforms persistence and seasonal naive baselines across the full evaluation year and during the vortex window, with a full-year MAE of 721 MW compared to 1,937 MW for the PJM Day-Ahead benchmark (both comparisons in the retrospective ERA5 setting). However, the central finding is not the point forecast improvement but the probabilistic failure: despite 86.8% empirical coverage at the 90% nominal prediction interval level over the full year, QR-GBT coverage falls to 66.7% during the vortex window, 98% interval coverage falls to 84.7%, and at the peak hour, the observed load exceeded the post-rearrangement QR-GBT q99 estimate of 139,585 MW, leaving the event peak outside the nominal 98% prediction interval. The 19.8-percentage-point gap between full-year and vortex-window coverage at the 90% level is invisible to any evaluation framework reporting only annual aggregate metrics — and it is this gap, rather than the point forecast comparison, that is the paper's primary empirical contribution.
1.6 Contributions and Organization
This paper makes three contributions:
Contribution 1. We provide a retrospective evaluation of QR-GBT during the January 6–8, 2014 PJM Polar Vortex, demonstrating that a model achieving 86.8% full-year empirical PI coverage nonetheless fails to cover the event peak of 140,510.2 MW at the nominal 98% level, with the observed load exceeding the post-rearrangement q99 estimate of 139,585 MW. All results use same-hour ERA5 retrospective reanalysis weather input. This constitutes a quantified, documented case of probabilistic calibration failure under out-of-distribution extreme cold.
Contribution 2. We show that the 19.8-percentage-point gap between full-year and vortex-window empirical coverage at the 90% level is invisible to any evaluation protocol reporting only annual aggregate metrics, and that calibration degrades systematically across the high-load tail of the demand distribution. These findings argue for disaggregated evaluation on historical extreme windows as a necessary component of probabilistic load forecasting benchmarks in reliability-critical applications.
Contribution 3. Using publicly available PJM Data Miner 2 load records and ERA5 reanalysis weather inputs, we construct a transparent retrospective benchmark on a historically significant event against which future extreme-event probabilistic forecasting methods can be compared, with all data sources, model specifications, and evaluation windows stated explicitly.
The remainder of this paper is organized as follows. Section 2 describes the PJM load dataset, ERA5 weather inputs, preprocessing procedure, and the definition of the polar vortex case study window. Section 3 presents the quantile regression forecasting framework, model architecture, feature set, and benchmark models. Section 4 describes the experimental design and evaluation protocol. Section 5 reports results across both evaluation windows. Section 6 discusses the mechanisms underlying calibration failure and their implications for probabilistic forecasting evaluation practice. Section 7 states the study's limitations. Section 8 concludes.